\scanpage{75}
Si $\varphi,\psi$ \inserted{$\in S_i$ ($i\in\sigma$)} sont dans une $=$ classe \inserted{$i\in\sigma$} et $\varphi\ne\psi$, on a $\varphi\cap\psi=\varnothing$ i.e. $\varphi^{-1}\psi$ sans pt fixe (en effet, c'est une des deux perm. circ. de $s$) donc les \inserted{trois} $\varphi\in S_i$ forment une \uline{partition} de l'ens $\Delta$ de cardinal $9$. De $\overline{\unsure{m|=}}$ pour l'ens $S_j$ ($j\in\sigma$, $j\ne i$). On trouve ainsi une application
\[
\Delta\xrightarrow{\sim}\prod_{i\in\sigma}S_i
\]
je dis qu'elle est bijective. \deleted{\unsure{Il|Le}}\unsure{Cela|Il} revient au $=$ de dire qu'elle est injective, ou encore de dire qu'un él. $(\alpha,\beta)$ de $\Delta=s\times t$ est connu quand on connaît ses deux composantes dans $S_i\times S_j$ ($\sigma=\{i,j\}$). Choisissons un isom $s\simeq t$, d'où un isom $\sigma\simeq\{1,2\}$, \deleted{\unsure{$\Delta$|et}} \inserted{$\Delta\to\underset{\text{''}S_1}{c}\times\underset{\text{''}S_2}{d}$,} il revient au $=$ de dire qu'un couple $(\alpha,\beta)$ d'él. de $s$ est connu quand on connaît l'unique permutation circulaire $\rho_1$ telle que $\beta=\rho_1\alpha$ (ceci détermine $\beta$ mod $\alpha$) puis l'unique transposition $\theta$ telle que $\beta=(\underbrace{\rho_1\theta}_{\rho_2})\alpha$ i.e. telle que $\theta\alpha=\alpha$ -- or $\alpha$ est \inserted{alors} le \inserted{unique} pt fixe de $\theta$.

\deleted{Considérons alors} Il est clair par construction que deux pts de $\Delta$ ayant $\overline{\unsure{m|=}}$ image dans $c$ ou dans $d$ sont \uline{non liés} dans $\Delta$, donc deux pts de $\Delta$ qui sont liés \deleted{\unsure{$\Delta$|a}} ont des images distinctes dans $c$ \uline{et} dans $d$, donc sont liés pour la \deleted{\unsure{relation|valeur}} \inserted{structure} de bloc \uline{de triades} \inserted{$D$} définie sur $\Delta$ grâce à la bijection $\Delta\simeq\prod_{i\in\sigma}S_i$. Comme les cardinal de l'ens des arêtes de $\Delta$ et de $D$ ont $=$ cardinal $18$, on conclut que les deux structures coïncident cqfd.

\uline{Conclusion} Dans la catégorie \inserted{(groupoïde \unsure{même|mieux})} $(\mathrm{Bitr})$ des bitriangles, le foncteur cov. $\Delta\longmapsto\overline\Delta$ définit donc une \deleted{\unsure{auto|autre}} \inserted{auto}équivalence involutive. On peut l'interpréter aussi comme une autoéquivalence involutive dans la catégorie $\unsure{\widehat{\mathrm{Ens}}|\mathrm{Ens}}_{\{3,3\}}$ des ens \inserted{$T$}. de cardinal $6$ munis d'une \deleted{\unsure{autoéquiv.}} \inserted{partition} de type $(3,3)$. Pour présenter de façon symétrique cette relation entre $T$ et $S$ de $\widehat{\mathrm{Ens}}_{(3,3)}$ on considère comme ``dualité'' entre $T$ et $S$
% archive sheet number: 74
